\begin{textsample}{NEG Dim 6 – vem\_ed\_35 – Score -22.00 – vem\_ed\_35/t188.txt}  \label{ex:f6_neg_007}
BOAS PRÁTICAS Com sabor de café e \textbf{chocolate} Alessandro José Padin Ferreira Fatec Praia Grande Nos dois primeiros semestres de 2025, vivi uma das experiências mais significativas da minha trajetória docente.
Por meio dos Projetos Colaborativos Internacionais (PCIs), desenvolvidos na abordagem COIL (Collaborative Online International Learning), tive a oportunidade de orientar dois PCIs entre a Fatec Praia Grande e o Departamento Universitario Obrero y Campesino (DUOC), do Chile.
Foram encontros virtuais que se transformaram em espaços de troca humana, cultural e \textbf{afetiva}.
No primeiro PCI, trabalhei com estudantes do curso de Comércio Exterior.
A proposta parecia objetiva: discutir \textbf{possibilidades} de inserção e fortalecimento do café brasileiro especial no mercado chileno, mas logo percebi que não seria possível falar apenas da mercadoria.
No Brasil, café também carrega memória, identidade e emoção.
Ao longo das atividades, incentivei os alunos a \textbf{compreenderem} que internacionalizar um produto também significa internacionalizar narrativas.
Não bastava apresentar números de exportação ou tendências de mercado.
Era necessário mostrar aos colegas chilenos como esse fruto faz parte da nossa história cotidiana.
Falamos sobre Santos, tão próxima da Fatec Praia Grande, cidade cuja trajetória está profundamente \textbf{ligada} ao porto por onde o produto brasileiro alcançou o mundo.
Conversamos sobre o cheiro do café passado nas manhãs brasileiras, sobre as reuniões em família e o \textbf{hábito} quase ritualístico de oferecer uma xícara a quem chega.
Em muitos momentos, percebi que os estudantes começaram a enxergar o comércio exterior de outra maneira.
Não apenas como \textbf{negociações} internacionais, mas como encontros culturais.
As apresentações produzidas por eles mostravam o crescimento do café brasileiro no Chile e a \textbf{valorização} \textbf{crescente} das alternativas especiais.
Mas o que mais me marcou foi observar como os alunos passaram a construir conexões emocionais em \textbf{torno} do tema.
Uma das equipes trouxe poemas e suas memórias \textbf{afetivas}.
Naquele \textbf{instante}, compreendi que o projeto havia alcançado algo maior do que eu imaginava inicialmente, pois deixava de ser apenas um produto de exportação e se tornava linguagem cultural.
Os estudantes chilenos demonstravam enorme curiosidade ao perceber essa \textbf{relação} \textbf{afetiva} do brasileiro.
Muitos comentavam como o Chile, historicamente \textbf{ligado} ao consumo de chá, vinha desenvolvendo novos \textbf{hábitos} \textbf{ligados} aos cafés especiais e às cafeterias contemporâneas.
As conversas eram \textbf{espontâneas}, e pouco a pouco as barreiras linguísticas iam desaparecendo. \textbf{Chocolate} \textbf{orgânico} No segundo semestre de 2025, \textbf{desenvolvi} um novo PCI, desta vez com estudantes de Gestão Empresarial.
O tema escolhido foi o \textbf{chocolate} \textbf{orgânico} brasileiro.
Mais uma vez, \textbf{procurei} estimular os alunos a \textbf{compreenderem} que por trás de um produto \textbf{existem} \textbf{territórios}, pessoas, tradições e modos de vida.
As discussões nos levaram ao sul da Bahia, às cooperativas \textbf{agrícolas}, à agricultura familiar e ao sistema cabruca, técnica tradicional de cultivo do cacau sob a sombra da Mata Atlântica.
Falamos sobre sustentabilidade, preservação ambiental e sobre como pequenas comunidades brasileiras transformaram o \textbf{chocolate} \textbf{orgânico} em símbolo de identidade regional.
Foi especialmente interessante observar o envolvimento dos estudantes ao pesquisarem o \textbf{universo} dos \textbf{chocolates} bean-to-bar (significa do “grão à barra” e define a produção de forma artesanal e integral), dos produtos de origem única e das marcas brasileiras que começam a conquistar mercados internacionais a partir de narrativas \textbf{ligadas} à sustentabilidade e à rastreabilidade.
Ao longo dos dois projetos, percebi algo que muitas vezes a rotina acadêmica nos faz esquecer: a educação possui um enorme poder de aproximação humana.
Mesmo separados por \textbf{milhares} de quilômetros, estudantes brasileiros e chilenos criaram vínculos \textbf{genuínos}.
No início, havia insegurança com o idioma, receio de falar espanhol, medo de errar.
Depois, surgiram risadas, espontaneidade e confiança.
Como professor, foi extremamente \textbf{gratificante} \textbf{acompanhar} esse processo.
Vi alunos descobrirem que a internacionalização não pertence apenas às grandes universidades ou aos programas presenciais no exterior.
Pode acontecer, também, dentro da sala de aula virtual, por meio do diálogo, da escuta e da colaboração.
Competências interculturais O reconhecimento institucional do projeto destacou competências interculturais, comunicação, habilidades digitais e trabalho em equipe global.
Mas acredito que houve um aprendizado ainda mais profundo: a compreensão de que produtos carregam histórias e que nenhuma estratégia internacional é realmente \textbf{eficiente} quando \textbf{ignora} a cultura das pessoas envolvidas.
Ao final daqueles encontros, tive a sensação de que café e \textbf{chocolate} eram apenas o ponto de partida.
O que realmente \textbf{atravessou} as fronteiras entre Brasil e Chile foram memórias, identidades e experiências compartilhadas.
E talvez seja essa a principal lembrança que guardarei dos PCIs de 2025: a certeza de que ensinar também é criar pontes humanas, algumas delas feitas de tecnologia, outras de palavras, e muitas delas do aroma inesquecível de um café recém passado.

% matched lemmas: acompanhar, afetivo, agrícola, atravessar, chocolate, compreender, crescente, desenvolvi, eficiente, espontâneo, existir, genuíno, gratificante, hábito, ignorar, instante, ligar, milhar, negociação, orgânico, possibilidade, procur, relação, território, torno, universo, valorização
\end{textsample}
